\textbf{人教B版选必1 第1章 空间向量与立体几何·讲练件（140题）}

全件140题：简单21｜中档104｜难15

{\def\LTcaptype{none} % do not increment counter
\begin{longtable}[]{@{}
  >{\raggedright\arraybackslash}p{(\linewidth - 6\tabcolsep) * \real{0.5234}}
  >{\raggedright\arraybackslash}p{(\linewidth - 6\tabcolsep) * \real{0.0721}}
  >{\raggedright\arraybackslash}p{(\linewidth - 6\tabcolsep) * \real{0.2853}}
  >{\raggedright\arraybackslash}p{(\linewidth - 6\tabcolsep) * \real{0.1191}}@{}}
\toprule\noalign{}
\endhead
\bottomrule\noalign{}
\endlastfoot
\textbf{节名} & \textbf{题量} & \textbf{简单/中档/难} &
\textbf{题型组数} \\
1.1.1 空间向量及其运算 & 10 & 简单6/中档4/难0 & 9 \\
1.1.2 空间向量基本定理 & 4 & 简单2/中档2/难0 & 3 \\
1.1.3 空间向量的坐标与空间直角坐标系 & 10 & 简单3/中档7/难0 & 7 \\
1.2.1 空间中的点、直线与空间向量 & 9 & 简单2/中档6/难1 & 6 \\
1.2.2 空间中的平面与空间向量 & 7 & 简单2/中档5/难0 & 5 \\
1.2.3 直线与平面的夹角 & 8 & 简单1/中档6/难1 & 7 \\
1.2.4 二面角 & 13 & 简单0/中档13/难0 & 12 \\
1.2.5 空间中的距离 & 79 & 简单5/中档61/难13 & 69 \\
合计 & 140 & 简单21/中档104/难15 & 118 \\
\end{longtable}
}

1.1 空间向量及其运算

\subsubsection{\texorpdfstring{\textbf{1.1
空间向量及其运算}}{1.1 空间向量及其运算}}\label{ux7a7aux95f4ux5411ux91cfux53caux5176ux8fd0ux7b97}

1.1.1 空间向量及其运算

\subsubsection{\texorpdfstring{\textbf{1.1.1
空间向量及其运算}}{1.1.1 空间向量及其运算}}\label{ux7a7aux95f4ux5411ux91cfux53caux5176ux8fd0ux7b97-1}

本节10题：简单6｜中档4｜难0　提分线：简单→保60\%｜中档→保80\%

\textbf{1.1.1.1 知识讲解｜空间向量及其运算}

1.1.1-1．〔基〕空间向量

【⁠编⁠注⁠】与必修2平面向量初步的定义类比记忆（只多``空间''二字）；零向量方向任意、模为0，是概念题易错点。

（1）定义：在空间，我们把具有大小和方向的量叫做空间向量．

（2）长度：空间向量的大小叫做空间向量的长度或模．

【⁠编⁠注⁠】对比辨析------空间向量与平面向量（旧知识：必修2第6章平面向量初步）：

{\def\LTcaptype{none} % do not increment counter
\begin{longtable}[]{@{}
  >{\raggedright\arraybackslash}p{(\linewidth - 4\tabcolsep) * \real{0.0565}}
  >{\raggedright\arraybackslash}p{(\linewidth - 4\tabcolsep) * \real{0.0941}}
  >{\raggedright\arraybackslash}p{(\linewidth - 4\tabcolsep) * \real{0.3083}}@{}}
\toprule\noalign{}
\endhead
\bottomrule\noalign{}
\endlastfoot
对比项 & 旧知识：平面向量 & 新知识：空间向量 \\
研究范围 & 同一平面内（二维） & 空间中（三维），平面情形为其特例 \\
定义 & 具有大小和方向的量 & 具有大小和方向的量，含义一致 \\
加减、数乘与数量积 & 三角形／平行四边形法则与运算律 &
法则与运算律完全一致，直接迁移 \\
共线与共面 & 共线向量定理刻画平面内位置关系 &
共线定理（条目5）与共面向量定理（条目9）并用，刻画空间位置关系 \\
易混点 & ------ &
任意两个空间向量必共面，但``不共线''的两向量只张成一张平面，不能作空间基底（须三个不共面，条目10） \\
\end{longtable}
}

1.1.1-2．〔基〕\textbf{空间向量的表示}

【⁠编⁠注⁠】字母表示与有向线段表示贯穿全章书写，解题前先统一记号；模的记号与绝对值区分，避免混淆。

（1）字母表示法：用字母\(\overset{⃑}{a},\overset{⃑}{b},\overset{⃑}{c}, \cdot \cdot \cdot\)，表示．

（2）几何表示法：用有向线段表示，其长度表示空间向量的模．即若向量\emph{a}的起点是\emph{A}，终点是\emph{B}，则向量\(\overset{⃑}{a}\)也可以记作\(\overset{⃑}{AB}\)，其模记为\(\left| \overset{⃑}{a} \right|\)或\(\left| \overset{⃑}{AB} \right|\)．

1.1.1-3．〔基〕几类特殊向量

【⁠编⁠注⁠】概念辨析高频条目，判定按定义逐条核对；注意共线向量不要求在同一直线上（平行或重合均可），且零向量与任意向量平行（衔接条目5）。

{\def\LTcaptype{none} % do not increment counter
\begin{longtable}[]{@{}
  >{\raggedright\arraybackslash}p{(\linewidth - 2\tabcolsep) * \real{0.0760}}
  >{\raggedright\arraybackslash}p{(\linewidth - 2\tabcolsep) * \real{0.3801}}@{}}
\toprule\noalign{}
\begin{minipage}[b]{\linewidth}\raggedright
特殊向量
\end{minipage} & \begin{minipage}[b]{\linewidth}\raggedright
定义
\end{minipage} \\
\midrule\noalign{}
\endhead
\bottomrule\noalign{}
\endlastfoot
零向量 & 长度为0的向量叫做零向量，记为0 \\
单位向量 & 模为1的向量叫做单位向量 \\
相反向量 &
与\(\overset{⃑}{a}\)长度相同而方向相反的向量叫做\(\overset{⃑}{a}\)的相反向量，记为\(- \overset{⃑}{a}\) \\
共线向量或平行向量 &
若表示若干空间向量的有向线段所在的直线互相平行或重合，那么这些向量叫做共线向量或平行向量；

规定：零向量与任意向量平行．即对任意向量\(\overset{⃑}{a}\)，都有\(\overset{⃑}{0} \parallel \overset{⃑}{a}\) \\
相等向量 & 方向相同且模相等的向量 \\
\end{longtable}
}

1.1.1-4．〔基〕空间向量的线性运算

【⁠编⁠注⁠】全章运算的基础，加减、数乘的法则与运算律均与平面向量一致，类比迁移即可；线性组合的运用见条目9与10。

空间向量的线性运算

{\def\LTcaptype{none} % do not increment counter
\begin{longtable}[]{@{}
  >{\raggedright\arraybackslash}p{(\linewidth - 4\tabcolsep) * \real{0.1178}}
  >{\raggedright\arraybackslash}p{(\linewidth - 4\tabcolsep) * \real{0.1804}}
  >{\raggedright\arraybackslash}p{(\linewidth - 4\tabcolsep) * \real{0.1583}}@{}}
\toprule\noalign{}
\begin{minipage}[b]{\linewidth}\raggedright
加法
\end{minipage} & \begin{minipage}[b]{\linewidth}\raggedright
\(\overset{⃑}{a} + \overset{⃑}{b} = \overset{⃑}{OA} + \overset{⃑}{AB} =\)\(\overset{⃑}{OB}\)
\end{minipage} &
\multirow{2}{=}{\begin{minipage}[b]{\linewidth}\raggedright
\includegraphics[width=0.95243in,height=0.4731in,alt={@@@57ceb128-6ced-4fe2-9e22-e16addd5170c}]{media/media/sub3_B_1.png}
\end{minipage}} \\
\begin{minipage}[b]{\linewidth}\raggedright
减法
\end{minipage} & \begin{minipage}[b]{\linewidth}\raggedright
\(\overset{⃑}{a} - \overset{⃑}{b} = \overset{⃑}{OA} - \overset{⃑}{OC} =\)\(\overset{⃑}{CA}\)
\end{minipage} \\
\midrule\noalign{}
\endhead
\bottomrule\noalign{}
\endlastfoot
\multirow{3}{=}{数乘运算} &
当\(\lambda > 0\)时，\(\lambda\overset{⃑}{a} = \lambda\overset{⃑}{OA} = \overset{⃑}{PQ}\)
&
\multirow{3}{=}{\includegraphics[width=0.85275in,height=0.7166in,alt={@@@e72cf1b0-1873-4d41-a581-224c31cd9465}]{media/media/sub3_B_2.png}} \\
&
当\(\lambda < 0\)时，\(\lambda\overset{⃑}{a} = \lambda\overset{⃑}{OA} = \overset{⃑}{MN}\) \\
& 当\(\lambda = 0\)时，\(\lambda\overset{⃑}{a} = \overset{⃑}{0}\) \\
\end{longtable}
}

运算律

{\def\LTcaptype{none} % do not increment counter
\begin{longtable}[]{@{}
  >{\raggedright\arraybackslash}p{(\linewidth - 2\tabcolsep) * \real{0.0984}}
  >{\raggedright\arraybackslash}p{(\linewidth - 2\tabcolsep) * \real{0.3577}}@{}}
\toprule\noalign{}
\begin{minipage}[b]{\linewidth}\raggedright
交换律
\end{minipage} & \begin{minipage}[b]{\linewidth}\raggedright
\(\overset{⃑}{a} + \overset{⃑}{b} =\)\(\overset{⃑}{b} + \overset{⃑}{a}\)
\end{minipage} \\
\midrule\noalign{}
\endhead
\bottomrule\noalign{}
\endlastfoot
结合律 &
\(\left( \overset{⃑}{a} + \overset{⃑}{b} \right) + \overset{⃑}{c} =\)\(\overset{⃑}{a} + \left( \overset{⃑}{b} + \overset{⃑}{c} \right)\)，\(\lambda(\mu\overset{⃑}{a}) =\)\((\lambda\mu)\overset{⃑}{a}\) \\
分配律 &
\((\lambda + \mu)\overset{⃑}{a} = \lambda\overset{⃑}{a} + \mu\overset{⃑}{a},\lambda(\overset{⃑}{a} + \overset{⃑}{b}) =\)\(\lambda\overset{⃑}{a} + \lambda\overset{⃑}{b}\) \\
\end{longtable}
}

1.1.1-5．〔基〕空间向量共线的充要条件

【⁠编⁠注⁠】证线线平行与求参数存在性的基础工具；易错点：与条目9共面充要条件混淆------共线用一个实数、共面用有序实数对。

对任意两个空间向量\(\overset{⃑}{a},\overset{⃑}{b}(\overset{⃑}{b} \neq 0)\)，\(\overset{⃑}{a} \parallel \overset{⃑}{b}\)的充要条件是存在实数\(\lambda\)，使\(\overset{⃑}{a} = \lambda\overset{⃑}{b}\)．

1.1.1-6．〔基〕空间向量的夹角

【⁠编⁠注⁠】夹角范围（0°到180°）是数量积正负讨论的前提；易混点：几何角（如异面直线所成角）与向量夹角可能互补，运算前先分清求的是哪个角。

{\def\LTcaptype{none} % do not increment counter
\begin{longtable}[]{@{}
  >{\raggedright\arraybackslash}p{(\linewidth - 2\tabcolsep) * \real{0.0510}}
  >{\raggedright\arraybackslash}p{(\linewidth - 2\tabcolsep) * \real{0.4051}}@{}}
\toprule\noalign{}
\begin{minipage}[b]{\linewidth}\raggedright
图示
\end{minipage} & \begin{minipage}[b]{\linewidth}\raggedright
\includegraphics[width=2.68733in,height=0.91736in,alt={@@@93d04d71-6a01-4600-baaf-bcb9db49ca41}]{media/media/sub3_B_3.png}
\end{minipage} \\
\midrule\noalign{}
\endhead
\bottomrule\noalign{}
\endlastfoot
定义 &
已知两个非零向量\(\overset{⃑}{a}\)，\(\overset{⃑}{b}\)，在空间任取一点\emph{O}，作\(\overset{⃑}{OA} = \overset{⃑}{a}\)，\(\overset{⃑}{OB} = \overset{⃑}{b}\)，则\(\angle AOB\)叫做向量\(\overset{⃑}{a}\)，\(\overset{⃑}{b}\)的夹角，记作\(\left\langle \overset{⃑}{a},\overset{⃑}{b} \right\rangle\) \\
范围 &
通常规定：0\(\leq \langle\overset{⃑}{a},\overset{⃑}{b}\rangle \leq\)\(\overset{⃑}{\pi}\)；

当\(\langle\overset{⃑}{a},\overset{⃑}{b}\rangle =\)\(\frac{\pi}{2}\)时，\(\overset{⃑}{a}\)与\(\overset{⃑}{b}\)垂直，记作\(\overset{⃑}{a}\bot\overset{⃑}{b}\) \\
\end{longtable}
}

1.1.1-7．〔基〕空间向量的数量积

【⁠编⁠注⁠】求空间角与距离的核心运算（支撑条目16、17、34）；注意数量积不满足结合律，零向量与任意向量的数量积为0。

（1）定义：已知两个非零向量\(\overset{⃑}{a}\)，\(\overset{⃑}{b}\)，则向量\(\overset{⃑}{b}\)的模长与\(\overset{⃑}{a}\)在向量\(\overset{⃑}{b}\)方向上的投影的乘积叫做\(\overset{⃑}{a}\)，\(\overset{⃑}{b}\)的数量积，记作\(\overset{⃑}{a} \cdot \overset{⃑}{b}\)．即\(\overset{⃑}{a} \cdot \overset{⃑}{b} =\)\(|\overset{⃑}{a}||\overset{⃑}{b}|\cos\langle\overset{⃑}{a},\overset{⃑}{b}\rangle\)．

【⁠微⁠提⁠醒⁠】零向量与任意向量的数量积为0．

（2）由数量积的定义，可以得到：\(\overset{⃑}{a}\bot\overset{⃑}{b} \Leftrightarrow\)\(\overset{⃑}{a} \cdot \overset{⃑}{b} = 0\)；\(\overset{⃑}{a} \cdot \overset{⃑}{a} = |\overset{⃑}{a}||\overset{⃑}{a}|\cos\langle\overset{⃑}{a},\overset{⃑}{a}\rangle =\)\({\overset{⃑}{a}}^{2}\)．

1.1.1-8．〔基〕投影向量

【⁠编⁠注⁠】数量积几何意义的载体，条目33三余弦定理与条目38点面距的推导都基于投影；易错点：投影向量是向量、投影长度是数量。

（1）在空间，向量\(\overset{⃑}{a}\)向向量\(\overset{⃑}{b}\)投影：

如图①，先将它们平移到同一平面\(\alpha\)内，利用平面上向量的投影，得到与向量\(\overset{⃑}{b}\)共线的向量\(\overset{⃑}{c}\)，\(\overset{⃑}{c} =\)\(\left| \overset{⃑}{a} \right|\cos\left\langle \overset{⃑}{a},\overset{⃑}{b} \right\rangle \cdot \frac{\overset{⃑}{b}}{\left| \overset{⃑}{b} \right|}\)，称向量\(\overset{⃑}{c}\)为向量\(\overset{⃑}{a}\)在向量\(\overset{⃑}{b}\)上的投影向量．

（2）向量\(\overset{⃑}{a}\)在直线\emph{l}上的投影如图②．

\includegraphics[width=3.38583in,height=0.8977in,alt={@@@5da5e1ea-9d33-4579-a56f-804b4a72c894}]{media/media/sub3_B_4.png}

（3）向量\(\overset{⃑}{a}\)向平面\(\beta\)投影：

如图③，分别由向量\(\overset{⃑}{a}\)的起点\emph{A}和终点\emph{B}作平面\(\beta\)的垂线，垂足分别为\(A'\)，\(B'\)，得到向量\(\overset{⃑}{A'B'}\)，向量\(\overset{⃑}{A'B'}\)称为向量\(\overset{⃑}{a}\)在平面\(\beta\)上的投影向量．

1.1.1-9．〔基〕空间向量共面的充要条件

【⁠编⁠注⁠】证四点共面、线共面的标准工具；易错点：定理前提是p与a、b不共线，此前提不可漏，且实数对存在唯一。

（1）共面向量：平行于同一平面的向量，叫做共面向量．

（2）空间向量共面的充要条件：向量\(\overset{⃑}{p}\)与不共线向量\(\overset{⃑}{a}\)，\(\overset{⃑}{b}\)共面的充要条件是存在唯一的有序实数对\((x,y)\)，使\(\overset{⃑}{p} =\)\(\overset{⃑}{}x\overset{⃑}{a} + y\overset{⃑}{b}\)．

\textbf{1.1.1.2 空间向量的概念辨析}

1题：-1

【⁠编⁠注⁠】题型通式：题干围绕空间向量的模、方向、零向量、单位向量、相等向量与相反向量罗列若干说法要求辨误时，判归本题型。第一步逐条对照定义核验每个说法；再对存疑条目举反例（模相等而方向不同、单位向量方向各异等）；最后排除错误项定答案。

\textbf{1.1.1.2-1．（简单）}下列说法正确的是（~~~~）

A．若\(\left| \overset{⃑}{a} \right| = \left| \overset{⃑}{b} \right|\)，则\(\overset{⃑}{a} = \overset{⃑}{b}\)或\(\overset{⃑}{a} = - \overset{⃑}{b}\)；B．若\(\overset{⃑}{a},\overset{⃑}{b}\)为相反向量，则\(\overset{⃑}{a} + \overset{⃑}{b} = \overset{⃑}{0}\)
；C．零向量是没有方向的向量；D．若\(\overset{⃑}{a},\overset{⃑}{b}\)是两个单位向量，则\(\overset{⃑}{a} = \overset{⃑}{b}\)

【⁠答⁠案⁠】

B　

【⁠知⁠识⁠点⁠】

1.1.1 向量的模、零向量与单位向量、相等向量、相反向量

【⁠分⁠析⁠】根据零向量、单位向量、向量的模、相等向量、相反向量等概念逐一分析即可得出答案.

【⁠详⁠解⁠】解：若\(\left| \overset{⃑}{a} \right| = \left| \overset{⃑}{b} \right|\)，则它们的方向相同时是相等向量，方向相反时是相反向量，还有可能方向既不相同，也不相反，A错；

若\(\overset{⃑}{a},\overset{⃑}{b}\)为相反向量，则它们的和为零向量，B对；

零向量的方向是任意的，C错；

两个单位向量只是模都为1，方向不一定相同，D错.故选：B.

\textbf{1.1.1.3 空间向量的线性运算}

1题：-2

【⁠编⁠注⁠】题型通式：题干在几何体中给出若干已知向量、要求用它们表示指定向量时，判归本题型。第一步为目标向量选一条首尾相接的路径写出加减链；再把链中各段用相等或共线的已知向量替换；最后合并化简得表达式。

\textbf{1.1.1.3-2．（简单）}在直三棱柱\(ABC - A_{1}B_{1}C_{1}\)中，若\(\overrightarrow{CA} = \overrightarrow{a},\overrightarrow{CB} = \overrightarrow{b},\overrightarrow{CC_{1}} = \overrightarrow{c},\)，则\(\overrightarrow{A_{1}B}\)=\_\_\_\_\_\_\_\_\_\_\_\_.(用\(\overrightarrow{a},\overrightarrow{b},\overrightarrow{c}\)表示)

【⁠答⁠案⁠】\(\overrightarrow{b} - \overrightarrow{a} - \overrightarrow{c}\)

【⁠知⁠识⁠点⁠】

1.1.1 棱柱的结构特征和分类、空间向量加减运算的几何表示

【⁠分⁠析⁠】连接\(CA_{1},\)根据直三棱柱的结构特征及空间向量减法的几何意义可得\(\overrightarrow{A_{1}B} = \overrightarrow{CB} - \overrightarrow{CA} - \overrightarrow{CC_{1}}\)，结合已知即可求表达式.

【⁠详⁠解⁠】连接\(CA_{1},\)则\(\overrightarrow{A_{1}B} = \overrightarrow{CB} - \overrightarrow{CA_{1}} = \overrightarrow{CB} - \overrightarrow{CA} - \overrightarrow{CC_{1}} =\)
\(\overrightarrow{b} - \overrightarrow{a} - \overrightarrow{c}\).

\includegraphics[width=1.61764in,height=2.59547in,alt={@@@4e3aa752498e415b9cc3d6d60f7bb3fe}]{media/media/image1.png}

故答案为：\(\overrightarrow{b} - \overrightarrow{a} - \overrightarrow{c}\)

\textbf{1.1.1.4 用基底表示向量}

1题：-3

【⁠编⁠注⁠】题型通式：题干指定三个不共面的向量作基底、要求把目标向量写成基底的线性组合并求系数时，判归本题型。第一步把目标向量沿几何路径分解到基向量方向；再把各段按比例关系代入基底；最后比较等式两边系数求出待定值。

\textbf{1.1.1.4-3．（简单）}已知正方体\emph{ABCD}－\emph{A\textsubscript{1}B\textsubscript{1}C\textsubscript{1}D\textsubscript{1}}中，\(\overset{⃑}{A_{1}E}\)＝\(\frac{1}{4}\)
\(\overset{⃑}{A_{1}C_{1}}\)，若\(\overset{⃑}{AE} =\)\emph{x}\(\overset{⃑}{AA_{1}}\)＋\emph{y}(\(\overset{⃑}{AB}\)＋\(\overset{⃑}{AD}\))，则\emph{x}＝\_\_\_\_\_\_\_\_，\emph{y}＝\_\_\_\_\_\_\_\_.

\includegraphics[width=1.65275in,height=1.65275in,alt={@@@b0b55c9f-4634-4bba-b096-bc6dd0cf7747}]{media/media/image2.png}

【⁠答⁠案⁠】

1 \(\frac{1}{4}\)　

【⁠知⁠识⁠点⁠】

1.1.1 空间向量的加减运算、空间向量的数乘运算、用空间基底表示向量

【⁠分⁠析⁠】结合空间向量的线性运算列方程，由此求得\(x,y\)的值.

【⁠详⁠解⁠】\(\overset{⃑}{AE} = \overset{⃑}{AA_{1}} + \overset{⃑}{A_{1}E} = \overset{⃑}{AA_{1}} + \frac{1}{4}\overset{⃑}{A_{1}C_{1}} = \overset{⃑}{AA_{1}} + \frac{1}{4}\left( \overset{⃑}{AB} + \overset{⃑}{AD} \right)\)，所以\(x = 1,y = \frac{1}{4}\).故答案为：\(1\)；\(\frac{1}{4}\)

\textbf{1.1.1.5 共面向量的判定}

1题：-4

【⁠编⁠注⁠】题型通式：题干给出 p=xa+yb 或 MP=xMA+yMB
一类线性表示与"共面"结论的互推说法判断时，判归本题型。第一步核对充要条件的前提（两向量不共线或三点不共线）是否齐备；再对缺前提的说法举共线反例排除；最后汇总正确序号定选项。

\textbf{1.1.1.5-4．（简单）}有下列说法:

①若\(\overset{⃑}{p} = x\overset{⃑}{a} + y\overset{⃑}{b}\)，则\(\overset{⃑}{p}\)与\(\overset{⃑}{a}\)，\(\overset{⃑}{b}\)共面；②若\(\overset{⃑}{p}\)与\(\overset{⃑}{a}\)，\(\overset{⃑}{b}\)共面，则\(\overset{⃑}{p}\)=\emph{x}\(\overset{⃑}{a}\)+\emph{y}\(\overset{⃑}{b}\)；③若\(\overset{⃑}{MP}\)=\emph{x}\(\overset{⃑}{MA}\)+\emph{y}\(\overset{⃑}{MB}\)，则\emph{P}，\emph{M}，\emph{A}，\emph{B}共面；④若\emph{P}，\emph{M}，\emph{A}，\emph{B}共面，则\(\overset{⃑}{MP}\)=\emph{x}\(\overset{⃑}{MA}\)+\emph{y}\(\overset{⃑}{MB}\).其中正确的是（~~~~）

A．①②③④；B．①③④；C．①③；D．②④

【⁠答⁠案⁠】

C　

【⁠知⁠识⁠点⁠】

1.1.1 判定空间向量共面

【⁠分⁠析⁠】利用空间向量共面定理逐一判断即可.

【⁠详⁠解⁠】若\(\overset{⃑}{a}\)，\(\overset{⃑}{b}\)共线，由\(\overset{⃑}{p}\)=\emph{x}\(\overset{⃑}{a}\)+\emph{y}\(\overset{⃑}{b}\)知\(\overset{⃑}{p}\)一定与\(\overset{⃑}{a}\)，\(\overset{⃑}{b}\)共面，

若\(\overset{⃑}{a}\)，\(\overset{⃑}{b}\)不共线，则满足共面定理，\(\overset{⃑}{p}\)与\(\overset{⃑}{a}\)，\(\overset{⃑}{b}\)共面，①对；

同理③对；若\(\overset{⃑}{p}\)与\(\overset{⃑}{a}\)，\(\overset{⃑}{b}\)共面，且\(\overset{⃑}{a}\)，\(\overset{⃑}{b}\)共线，则不一定有\(\overset{⃑}{p}\)=\emph{x}\(\overset{⃑}{a}\)+\emph{y}\(\overset{⃑}{b}\)，故②不对；同理④不对，故选：C.

\textbf{1.1.1.6 数量积的概念与运算律}

1题：-5

【⁠编⁠注⁠】题型通式：题干列出数量积相关的若干等式（平方、商式、乘积展开）要求判断正误时，判归本题型。第一步按定义
a·b=\textbar a\textbar\textbar b\textbar cosθ
逐式展开核验；再警惕向量无除法、数量积无结合律等误用；最后汇总正确式定选项。

\textbf{1.1.1.6-5．（简单）}设\(\overset{⃑}{a}\)，\(\overset{⃑}{b}\)为空间中的任意两个非零向量，下列各式中正确的有（~~~~）

A．\({\overset{⃑}{a}}^{2} = \left| \overset{⃑}{a} \right|^{2}\)；B．\(\frac{\overset{⃑}{a} \cdot \overset{⃑}{b}}{\overset{⃑}{a} \cdot \overset{⃑}{a}} = \frac{\overset{⃑}{b}}{\overset{⃑}{a}}\)；C．\(\left( \overset{⃑}{a} \cdot \overset{⃑}{b} \right)^{2} = {\overset{⃑}{a}}^{2} \cdot {\overset{⃑}{b}}^{2}\)；D．\(\left( \overset{⃑}{a} - \overset{⃑}{b} \right)^{2} = {\overset{⃑}{a}}^{2} - 2\overset{⃑}{a} \cdot \overset{⃑}{b} + {\overset{⃑}{b}}^{2}\)

【⁠答⁠案⁠】

AD　

【⁠知⁠识⁠点⁠】

1.1.1 空间向量数量积的概念辨析、求空间向量的数量积

【⁠分⁠析⁠】根据空间向量数量积的定义与运算律一一判断即可；

【⁠详⁠解⁠】解：对于A：\({\overset{⃑}{a}}^{2} = \overset{⃑}{a} \cdot \overset{⃑}{a} = \left| \overset{⃑}{a} \right| \cdot \left| \overset{⃑}{a} \right|\cos 0 = \left| \overset{⃑}{a} \right|^{2}\)，故A正确；

对于B：因为向量不能做除法，即\(\frac{\overset{⃑}{b}}{\overset{⃑}{a}}\)无意义，故B错误；

对于C：\(\left( \overset{⃑}{a} \cdot \overset{⃑}{b} \right)^{2} = \left( \left| \overset{⃑}{a} \right| \cdot \left| \overset{⃑}{b} \right|\cos\left\langle \overset{⃑}{a},\overset{⃑}{b} \right\rangle \right)^{2} = \left| \overset{⃑}{a} \right|^{2} \cdot \left| \overset{⃑}{b} \right|^{2}\cos^{2}\left\langle \overset{⃑}{a},\overset{⃑}{b} \right\rangle\)，故C错误；

对于D：\(\left( \overset{⃑}{a} - \overset{⃑}{b} \right)^{2} = \left( \overset{⃑}{a} - \overset{⃑}{b} \right) \cdot \left( \overset{⃑}{a} - \overset{⃑}{b} \right) = {\overset{⃑}{a}}^{2} - 2\overset{⃑}{a} \cdot \overset{⃑}{b} + {\overset{⃑}{b}}^{2}\)，故D正确；故选：AD

\textbf{1.1.1.7 数量积求夹角与投影（非坐标法）}

1题：-6

【⁠编⁠注⁠】题型通式：题干在几何体中求两向量的夹角、模长或投影而不建坐标系时，判归本题型。第一步平移向量使共起点或选不共面的三个向量作基底；再把相关向量用基底表示并算出数量积与模长；最后代入夹角公式并按向量夹角范围
{[}0,π{]} 取角。

\textbf{1.1.1.7-6．（中档）}如图，在正方体\(ABCD - A_{1}B_{1}C_{1}D_{1}\)中，求向量\(\overset{⃑}{BC_{1}}\)与\(\overset{⃑}{AC}\)的夹角的大小．

\includegraphics[width=3.38583in,height=3.22956in,alt={@@@d8ffc038-cfd6-4d16-b846-f3d81f29b325}]{media/media/image3.png}

【⁠答⁠案⁠】\(60{^\circ}\)

【⁠知⁠识⁠点⁠】

1.1.1 空间向量数量积的应用

【⁠分⁠析⁠】方法1：结合图形，平移向量，利用空间向量的夹角定义来求，但要注意向量夹角的范围；

方法2：先求\(\overset{⃑}{BC_{1}} \cdot \overset{⃑}{AC}\)，再利用公式\(\cos\left\langle \overset{⃑}{BC_{1}},\overset{⃑}{AC} \right\rangle = \frac{\overset{⃑}{BC_{1}} \cdot \overset{⃑}{AC}}{\left| \overset{⃑}{BC_{1}} \right| \cdot \left| \overset{⃑}{AC} \right|}\)求\(\cos\left\langle \overset{⃑}{BC_{1}},\overset{⃑}{AC} \right\rangle\)，最后确定\(\left\langle \overset{⃑}{BC_{1}},\overset{⃑}{AC} \right\rangle\)即可．

【⁠详⁠解⁠】解：方法1：因为\(\overset{⃑}{AD_{1}} = \overset{⃑}{BC_{1}}\)，所以\(\angle CAD_{1}\)的大小就等于\(\left\langle \overset{⃑}{BC_{1}},\overset{⃑}{AC} \right\rangle\)

因为△\(CAD_{1}\)为等边三角形，所以\(\angle CAD_{1} = 60^{{^\circ}}\)，所以\(\overset{⃑}{BC_{1}}\)与\(\overset{⃑}{AC}\)的夹角的大小为\(60{^\circ}\)．

方法2．设正方体的棱长为1，\(\overset{⃑}{BC_{1}} \cdot \overset{⃑}{AC} = \left( \overset{⃑}{BC} + \overset{⃑}{CC_{1}} \right) \cdot \left( \overset{⃑}{AB} + \overset{⃑}{BC} \right) = \left( \overset{⃑}{AD} + \overset{⃑}{AA_{1}} \right) \cdot \left( \overset{⃑}{AB} + \overset{⃑}{AD} \right)\)\(= \overset{⃑}{AD} \cdot \overset{⃑}{AB} + \left| \overset{⃑}{AD} \right|^{2} + \overset{⃑}{AA_{1}} \cdot \overset{⃑}{AB} + \overset{⃑}{AA_{1}} \cdot \overset{⃑}{AD} = 0 + \left| \overset{⃑}{AD} \right|^{2} + 0 + 0 = \left| \overset{⃑}{AD} \right|^{2} = 1\)又因为\(\left| \overset{⃑}{BC_{1}} \right| = \sqrt{2},\left| \overset{⃑}{AC} \right| = \sqrt{2}\)，所以\(\cos\left\langle \overset{⃑}{BC_{1}},\overset{⃑}{AC} \right\rangle = \frac{\overset{⃑}{BC_{1}} \cdot \overset{⃑}{AC}}{\left| \overset{⃑}{BC_{1}} \right| \cdot \left| \overset{⃑}{AC} \right|} = \frac{1}{\sqrt{2} \times \sqrt{2}} = \frac{1}{2}\)，
因为\(\left\langle \overset{⃑}{BC_{1}},\overset{⃑}{AC} \right\rangle \in \lbrack 0,\pi\rbrack\)，所以\(\overset{⃑}{BC_{1}}\)与\(\overset{⃑}{AC}\)的夹角的大小为\(60{^\circ}\)．

\textbf{1.1.1.8 数量积条件求参（夹角与垂直）}

2题：-7～8

【⁠编⁠注⁠】题型通式：题干给出两向量的模与夹角、并以夹角为锐角/钝角或两向量垂直为条件求参数时，判归本题型。第一步用数量积把条件翻译成不等式或方程（锐角\textgreater0、钝角\textless0、垂直=0）；再求解并剔除使两向量共线的参数值；最后写出参数值或取值范围。

\textbf{1.1.1.8-7．（中档）}已知向量\(\overset{⃑}{a},\overset{⃑}{b}\)，且\(\left| \overset{⃑}{a} \right| = 2\)，\(\left| \overset{⃑}{b} \right| = 1\)，\(\left\langle \overset{⃑}{a},\overset{⃑}{b} \right\rangle = 60^{\circ}\)，则使向量\(\overset{⃑}{a} + \lambda\overset{⃑}{b}\)与\(\lambda\overset{⃑}{a} - 2\overset{⃑}{b}\)的夹角为钝角的实数\(\lambda\)的取值范围是\_\_\_\_\_\_．

【⁠答⁠案⁠】\(\left( - 1 - \sqrt{3}, - 1 + \sqrt{3} \right)\)

【⁠知⁠识⁠点⁠】

1.1.1 已知向量共线（平行）求参数、用定义求向量的数量积、向量夹角的计算

【⁠分⁠析⁠】根据\(\left( \overset{⃑}{a} + \lambda\overset{⃑}{b} \right) \cdot \left( \lambda\overset{⃑}{a} - 2\overset{⃑}{b} \right) < 0\)，得出方程求得\(- 1 - \sqrt{3} < \lambda < - 1 + \sqrt{3}\)，若向量\(\overset{⃑}{a} + \lambda\overset{⃑}{b}\)与\(\lambda\overset{⃑}{a} - 2\overset{⃑}{b}\)反向共线时，则存在实数\(k < 0\)使得\(\overset{⃑}{a} + \lambda\overset{⃑}{b} = k(\lambda\overset{⃑}{a} - 2\overset{⃑}{b})\)成立，得到方程组\(\left\{ \begin{array}{r}
k\lambda = 1 \\
\lambda = - 2k
\end{array} \right.\ \)无解，即可得到答案.

【⁠详⁠解⁠】由向量\(\overset{⃑}{a},\overset{⃑}{b}\)，且\(\left| \overset{⃑}{a} \right| = 2\)，\(\left| \overset{⃑}{b} \right| = 1\)，\(\left\langle \overset{⃑}{a},\overset{⃑}{b} \right\rangle = 60^{\circ}\)，可得\(\overset{⃑}{a} \cdot \overset{⃑}{b} = 2 \times 1 \times \cos 60^{\circ} = 1\)，

因为向量\(\overset{⃑}{a} + \lambda\overset{⃑}{b}\)与\(\lambda\overset{⃑}{a} - 2\overset{⃑}{b}\)的夹角为钝角，可得\(\cos\left\langle \overset{⃑}{a} + \lambda\overset{⃑}{b},\lambda\overset{⃑}{a} - 2\overset{⃑}{b} \right\rangle < 0\)且\(\cos\left\langle \overset{⃑}{a} + \lambda\overset{⃑}{b},\lambda\overset{⃑}{a} - 2\overset{⃑}{b} \right\rangle \neq - 1\)，由\(\left( \overset{⃑}{a} + \lambda\overset{⃑}{b} \right) \cdot \left( \lambda\overset{⃑}{a} - 2\overset{⃑}{b} \right) < 0\)，即\(\lambda{\overset{⃑}{a}}^{2}\text{+(}\lambda^{2} - 2)\overset{⃑}{a} \cdot \overset{⃑}{b} - 2\lambda{\overset{⃑}{b}}^{2} < 0\)，所以\(\lambda^{2} + 2\lambda - 2 < 0\)，解得\(- 1 - \sqrt{3} < \lambda < - 1 + \sqrt{3}\)，

若向量\(\overset{⃑}{a} + \lambda\overset{⃑}{b}\)与\(\lambda\overset{⃑}{a} - 2\overset{⃑}{b}\)反向共线时，存在实数\(k < 0\)，使得\(\overset{⃑}{a} + \lambda\overset{⃑}{b} = k(\lambda\overset{⃑}{a} - 2\overset{⃑}{b})\)成立，可得\(\left\{ \begin{array}{r}
k\lambda = 1 \\
\lambda = - 2k
\end{array} \right.\ \)，此时方程组无解，所以不存在实数\(k\)使得向量\(\overset{⃑}{a} + \lambda\overset{⃑}{b}\)与\(\lambda\overset{⃑}{a} - 2\overset{⃑}{b}\)反向共线，所以实数\(\lambda\)的取值范围是\(\left( - 1 - \sqrt{3}, - 1 + \sqrt{3} \right)\).

\textbf{1.1.1.8-8．（简单）}已知\(\left| \overset{⃑}{a} \right| = 3\sqrt{2},\left| \overset{⃑}{b} \right| = 4,\overset{⃑}{m} = \overset{⃑}{a} + \overset{⃑}{b},\overset{⃑}{n} = \overset{⃑}{a} + \lambda\overset{⃑}{b},\left\langle \overset{⃑}{a},\overset{⃑}{b} \right\rangle = 135{^\circ},\overset{⃑}{m}\bot\overset{⃑}{n}\)
，则λ＝\_\_\_\_\_\_\_\_．

【⁠答⁠案⁠】\(- \frac{3}{2}\)

【⁠知⁠识⁠点⁠】

1.1.1 用定义求向量的数量积、垂直关系的向量表示

【⁠分⁠析⁠】根据题意\(\overset{⃑}{m}\bot\overset{⃑}{n}\)得到\(\left( \overset{⃑}{a} + \overset{⃑}{b} \right) \cdot \left( \overset{⃑}{a} + \lambda\overset{⃑}{b} \right) = 0\)，从而可求出\(\lambda\)的值.

【⁠详⁠解⁠】由\(\overset{⃑}{m}\bot\overset{⃑}{n}\)，得\(\left( \overset{⃑}{a} + \overset{⃑}{b} \right) \cdot \left( \overset{⃑}{a} + \lambda\overset{⃑}{b} \right) = 0\)，即\({\overset{⃑}{a}}^{2} + (\lambda + 1) \cdot \overset{⃑}{a} \cdot \overset{⃑}{b} + \lambda{\overset{⃑}{b}}^{2} = 0\)，即
\(\left| \overset{⃑}{a} \right|^{2} + (\lambda + 1) \cdot \left| \overset{⃑}{a} \right| \cdot \left| \overset{⃑}{b} \right| \cdot \cos\left\langle \overset{⃑}{a},\overset{⃑}{b} \right\rangle + \lambda\left| \overset{⃑}{b} \right|^{2} = 0\)，所以\(18 + (\lambda + 1) \times 3\sqrt{2} \times 4 \times \cos 135{^\circ} + 16\lambda = 0\)，即4\emph{λ}＋6＝0，∴\emph{λ}＝\(- \frac{3}{2}\)．故答案为：\(- \frac{3}{2}\).

\textbf{1.1.1.9 数量积求距离（展开法）}

1题：-9

【⁠编⁠注⁠】题型通式：题干在含二面角的拼接图形中求两点间距离时，判归本题型。第一步把待求线段对应的向量写成已知棱向量的首尾相接和链；再对模长平方作数量积展开，逐项代入模长与夹角（端点处两垂线段的夹角由二面角决定，注意取
45° 还是 135°）；最后开方得距离。

\textbf{1.1.1.9-9．（中档）}如图，在大小为45°的二面角A\-EF\-D中，四边形ABFE，CDEF都是边长为1的正方形，则B，D两点间的距离是()

\includegraphics[width=2.61024in,height=1.4087in,alt={@@@494a628b97d74c2894c135f69e8f1eb7}]{media/media/image4.png}

A．\(\sqrt{3}\)；B．\(\sqrt{2}\)；C．1；D．\(\sqrt{3 - \sqrt{2}}\)

【⁠答⁠案⁠】

D　

【⁠知⁠识⁠点⁠】

1.1.1 空间向量数量积的应用

【⁠分⁠析⁠】由\(\overset{⃑}{DB} = \overset{⃑}{ED} + \overset{⃑}{FE} + \overset{⃑}{BF}\)，利用数量积运算性质展开即可得到答案

【⁠详⁠解⁠】\(\because\overset{⃑}{BD} = \overset{⃑}{ED} + \overset{⃑}{FE} + \overset{⃑}{BF}\),\(\therefore\left| \overset{⃑}{BD} \right|^{2} = \left| \overset{⃑}{BF} \right|^{2} + \left| \overset{⃑}{FE} \right|^{2} + \left| \overset{⃑}{ED} \right|^{2} + 2\overset{⃑}{BF} \bullet \overset{⃑}{FE} + 2\overset{⃑}{FE} \bullet \overset{⃑}{ED} + 2\overset{⃑}{BF} \bullet \overset{⃑}{ED} = 1 + 1 + 1 - \sqrt{2}\)故\(\left| \overset{⃑}{BD} \right| = \sqrt{3 - \sqrt{2}}\)故选\(D\)

【⁠点⁠睛⁠】本题是要求空间两点之间的距离，运用空间向量将其表示，然后计算得到结果，较为基础．

\textbf{1.1.1.10 数量积求距离（折叠矩形）}

1题：-10

【⁠编⁠注⁠】题型通式：题干为平面图形沿一条直线折起且折后两面垂直、求两点间距离时，判归本题型。第一步在折后图形中过两点向折痕作垂线定出垂足；再把目标向量写成两段垂线与折痕段的和链并平方展开；最后利用面面垂直带来的向量垂直逐项化简、开方得距离。

\textbf{1.1.1.10-10．（中档）}已知矩形\(ABCD\)中，\(AB = 1\)，\(BC = \sqrt{3}\)，将矩形\(ABCD\)沿对角线\(AC\)折起，使平面\(ABC\)与平面\(ACD\)垂直，则\(|\overset{⃑}{BD}| =\)（~~~~）．

A．\(\frac{\sqrt{10}}{2}\)；B．\(\frac{\sqrt{6}}{2}\)；C．\(\frac{\sqrt{5}}{2}\)；D．2

【⁠答⁠案⁠】

A　

【⁠知⁠识⁠点⁠】

1.1.1 空间向量数量积的应用

【⁠分⁠析⁠】过点\(B\)，\(D\)分别向\(AC\)作垂线，垂足分别为\(M\)，\(N\)，由\(\overset{⃑}{BD} = \overset{⃑}{BM} + \overset{⃑}{MN} + \overset{⃑}{ND}\)，平方后结合长度和垂直关系可得解.

【⁠详⁠解⁠】
\includegraphics[width=2.15169in,height=1.60713in,alt={@@@e34e5148-bca3-440e-84b1-3cbf33fe2cc1}]{media/media/image5.png}~~过点\(B\)，\(D\)分别向\(AC\)作垂线，垂足分别为\(M\)，\(N\)，则可得\(AM = \frac{1}{2}\)，\(BM = \frac{\sqrt{3}}{2}\)，\(CN = \frac{1}{2}\)，\(DN = \frac{\sqrt{3}}{2}\)，\(MN = 1\)．由于\(\overset{⃑}{BD} = \overset{⃑}{BM} + \overset{⃑}{MN} + \overset{⃑}{ND}\)，所以\(|\overset{⃑}{BD}|^{2} = {(\overset{⃑}{BM} + \overset{⃑}{MN} + \overset{⃑}{ND})}^{2}\)\(= |\overset{⃑}{BM}|^{2} + |\overset{⃑}{MN}|^{2} + |\overset{⃑}{ND}|^{2} + 2(\overset{⃑}{BM} \cdot \overset{⃑}{MN} + \overset{⃑}{MN} \cdot \overset{⃑}{ND} + \overset{⃑}{BM} \cdot \overset{⃑}{ND})\)\(= \left( \frac{\sqrt{3}}{2} \right)^{2} + 1^{2} + \left( \frac{\sqrt{3}}{2} \right)^{2} + 2 \times (0 + 0 + 0) = \frac{5}{2}\)，所以\(|\overset{⃑}{BD}| = \frac{\sqrt{10}}{2}\)．故选:\emph{A}．

【⁠点⁠睛⁠】本题主要考查了空间向量数量积的应用，属于基础题.
